problem:  
Let \((M^{n+1},g_+)\) be a conformally compact Einstein manifold with boundary conformal class \([\!h\!]\) on \(\partial M\).  
For each smooth representative \(h\in[\!h\!]\), the scattering construction produces the **fractional conformally covariant boundary operator** \(P^h_1\) (the order‑1 conformal fractional Laplacian), which coincides with the normalized Dirichlet‑to‑Neumann operator of \(g_+\).  
Its spectral zeta function
\[
\zeta_{P^h_1}(s) = \mathrm{Tr}\big((P^h_1)^{-s}\big)
\]
admits meromorphic continuation, and one defines the **zeta‑regularized determinant**
\[
\mathcal D(h) = \exp\!\big(-\zeta'_{P^h_1}(0)\big).
\]

---

**Open problem (Explicit anomaly and extremality of the Steklov determinant in high dimensions):**

1. (**Explicit anomaly functional**)  
   For odd \(n\ge3\), it is known that \(\log \mathcal D(h)\) has a local conformal anomaly of the form  
   \[
   \delta_\varphi \log \mathcal D(h)
      = \int_{\partial M} Q^{(1)}_h\, \varphi\, dA_h,
   \]
   where \(Q^{(1)}_h\) is a scalar curvature invariant polynomial in the jets of \(h\).  
   The general existence of such a \(Q^{(1)}_h\) follows from determinant‑anomaly theory.  
   The **open question** is to determine an explicit closed‑form formula for \(Q^{(1)}_h\) in all dimensions \(n\ge5\):
   - Can one express \(Q^{(1)}_h\) purely in terms of intrinsic and extrinsic curvature tensors of \(h\)—for example, as a universal linear combination of Branson's \(Q_n(h)\), the norm‑square of the second fundamental form, and their covariant derivatives?  
   - Does it coincide (up to normalization) with the Graham–Juhl holographic \(Q\)-curvature associated to the Einstein filling?

2. (**Extremal conformal metrics**)  
   If the boundary area \(A(\partial M,h)\) is fixed, do conformal critical points of \(\log \mathcal D(h)\) within \([\!h\!]\) satisfy the PDE
   \[
   Q^{(1)}_h = \text{constant on }\partial M?
   \]
   Equivalently, are stationary metrics precisely those whose fractional self‑interaction potential—encoded by \(P^h_1\)—has constant anomaly density?

3. (**Rigidity**)  
   In the subclass of conformally flat boundaries, does the constancy of \(Q^{(1)}_h\) imply that \(h\) is conformally equivalent to the round metric on \(S^n\)?  
   (In dimension \(n=3\), such a statement would mirror the uniqueness of the round metric as the critical point of the standard Polyakov functional.)

---

Why is it a "good" problem:  
This formulation pinpoints the genuinely unknown frontier: explicit determination and geometric interpretation of the **conformal anomaly density \(Q^{(1)}_h\)** for the order‑1 fractional conformal Laplacian beyond low dimensions.  
The existence of such an anomaly is abstractly guaranteed, but its concrete expression and geometric meaning remain open except in the lowest dimensions. Solving this would require blending scattering theory, holographic renormalization, and conformal invariant theory.  

Moreover, connecting its constancy to extremal metrics would create a boundary counterpart of the classical Polyakov problem, linking determinant geometry, conformal PDEs, and curvature flows. The problem is technically precise, deeply tied to active research (fractional operators, holographic anomalies, zeta determinants), and conceptually elegant—an archetype of a “narrow entrance, profound depth’’ question at the interface of spectral and conformal geometry.